%------------------------------------------------------------------------------%
\begin{question}
  Uma força é representada pelo vetor
  \[
    \vec{F} = (6, -2, 3)
  \]
  e atua sobre um objeto que se desloca na direção
  \[
    \vec{d} = (2, 1, 2)
  \]
  Determine as componentes da força paralela e perpendicular à direção $\vec{d}$
\end{question}

%------------------------------------------------------------------------------%
\begin{resolution}{Solução}
\begin{columns}[t]
\column{0.5\textwidth}
  \begin{align*}
    \onslide<+->{\vec{F} \cdot \vec{d}}
    \onslide<+->{& = \vetortri{6}{-2}{3} \cdot \vetortri{2}{1}{2} \\}
    \onslide<+->{& = 6 \times 2 + (-2)1 + 3 \times 2 \\}
    \onslide<+->{& = 12 - 2 + 6\\}
    \onslide<+->{& = 16}
  \end{align*}
\column{0.5\textwidth}
  \begin{align*}
    \onslide<+->{\vec{d}\cdot\vec{d}}
    \onslide<+->{& = \vetortri{2}{1}{2} \cdot \vetortri{2}{1}{2} \\}
    \onslide<+->{& = 2^2 + 1^2 + 2^2 \\}
    \onslide<+->{& = 4 + 1 + 4 \\}
    \onslide<+->{& = 9}
  \end{align*}
\end{columns}
\end{resolution}

%------------------------------------------------------------------------------%
\begin{resolution}{Solução}
  \[
    \onslide<+->{\projuv{F}{d}}
    \onslide<+->{\;=\; \frac{\vec{F}\cdot\vec{d}}{\vec{d}\cdot\vec{d}} \, \vec{d}}
    \onslide<+->{\;=\; \frac{16}{9} \vetortri{2}{1}{2}}
    \onslide<+->{\;=\; \Vetortri{\frac{16\times 2}{9}}{\frac{16}{9}}{\frac{16 \times 2}{9}}}
    \onslide<+->{\;=\; \Vetortri{\frac{32}{9}}{\frac{16}{9}}{\frac{32}{9}}}
  \]
\end{resolution}

%------------------------------------------------------------------------------%
\begin{resolution}{Solução}
  A componente paralela da força é
  \[
    \vec{F}_{\parallel}
    \;=\; \projuv{F}{d}
    \;=\; \Vetortri{\frac{32}{9}}{\frac{16}{9}}{\frac{32}{9}}
  \]
\end{resolution}

%------------------------------------------------------------------------------%
\begin{resolution}{Solução}
  A componente perpendicular é
  \[
    \vec{F}_{\perp}
    \;=\; \vec{F} - \vec{F}_{\parallel}
    \;=\; \vetortri{6}{-2}{3} - \Vetortri{\frac{32}{9}}{\frac{16}{9}}{\frac{32}{9}}
    \;=\; \Vetortri{\frac{ 6 \times 9 - 32}{9}}
                   {\frac{-2 \times 9 - 16}{9}}
                   {\frac{ 3 \times 9 - 32}{9}}
    \;=\; \Vetortri{\frac{22}{9}}{-\frac{34}{9}}{-\frac{5}{9}}
  \]
\end{resolution}

%------------------------------------------------------------------------------%
\endinput
